\chapter{Stomatitis}

\section*{Definition}
Stomatitis is a broad term for inflammation of the oral mucosa, including the buccal surfaces, lips, palate, tongue, and floor of mouth, presenting as redness, swelling, ulceration, or vesicle formation with associated pain.

\section*{Pathophysiology}
Direct cytotoxicity (chemotherapy, radiation), immune-mediated epithelial destruction (aphthous, autoimmune), or infectious invasion (viral, fungal, bacterial) disrupts the mucosal barrier integrity. The resulting inflammation recruits neutrophils and lymphocytes, producing pain, swelling, and functional impairment of eating and speaking.

\section*{Causes}
\begin{itemize}
  \item \textbf{Aphthous stomatitis (canker sores):} unknown cause minor, major, or herpetiform ulcers; associated with stress, trauma, and food sensitivities
  \item \textbf{Infectious:} HSV-1 (herpetic gingivostomatitis), Coxsackie (herpangina, hand-foot-mouth), VZV, Candida albicans (atrophic or pseudomembranous), EBV
  \item \textbf{Chemotherapy/radiation:} Mucositis from cytotoxic agents (methotrexate, 5-FU, cisplatin) and head and neck radiation
  \item \textbf{Autoimmune:} Pemphigus vulgaris, oral lichen planus, redness multiforme, lupus erythematosus
  \item \textbf{Allergic:} Contact stomatitis from dental materials, toothpastes, mouthwashes, foods
  \item \textbf{Systemic disease:} Crohn disease, Behcet disease, celiac disease, HIV/AIDS
\end{itemize}

\section*{When to See a Doctor}
See your doctor if:
\begin{itemize}
  \item Oral ulcers lasting > 2 weeks without healing
  \item Inability to eat, drink, or swallow due to pain
  \item Systemic symptoms: fever, rash, conjunctivitis, genital ulcers
  \item Vesicles or bullae on the oral mucosa or skin
  \item Immunocompromised patient with oral lesions (risk of systemic dissemination)
  \item Rapidly progressive ulceration or tissue tissue death
\end{itemize}

\section*{Related Symptoms}
\begin{itemize}
  \item Oral mucosal redness, erosions, or ulcerations
  \item painful swallowing (painful swallowing) and difficulty swallowing
  \item Sialorrhea (drooling) or dry mouth
  \item Halitosis and dysgeusia
  \item Fever and cervical lymphadenopathy (if infectious)
\end{itemize}
\textbf{Important:} Chemotherapy-induced mucositis affects up to 40\% of people receiving standard-dose chemotherapy and nearly all people receiving head and neck radiation; it is a dose-limiting toxicity and a risk factor for sepsis.

\section*{Self-Care / Home Management}
\begin{itemize}
  \item Soft, bland, non-acidic diet (avoid spicy, salty, and abrasive foods)
  \item Oral hydration (frequent sips of water, ice chips, popsicles)
  \item Saline or sodium bicarbonate mouth rinses (1/2 tsp salt or baking soda in 8 oz water)
  \item Topical barrier agents: milk of magnesia, sucralfate suspension, or Orabase
  \item Avoid alcohol, tobacco, and commercial mouthwashes containing alcohol
  \item Pain management: topical viscous lidocaine 2\% (swish and spit) before meals
  \item Good oral hygiene with a soft-bristled toothbrush and non-irritating toothpaste
\end{itemize}

\section*{How Your Doctor Will Evaluate You}
\begin{itemize}
  \item Comprehensive oral examination: characterize lesions (ulcer, vesicle, bulla, erosion), distribution, and severity
  \item Viral culture, PCR, or direct fluorescent antibody for HSV, VZV, Coxsackie
  \item KOH preparation and fungal culture if candidiasis suspected
  \item HIV testing if risk factors present
  \item Biopsy with direct immunofluorescence for autoimmune blistering disorders (pemphigus, lichen planus)
  \item Patch testing for suspected contact allergy
  \item Endoscopy and colonoscopy if Crohn disease or Behcet suspected
\end{itemize}

\section*{Treatment Options}
\begin{itemize}
  \item \textbf{Aphthous:} Topical corticosteroids (triamcinolone dental paste, fluocinonide gel), amlexanox paste, or cautery (silver nitrate)
  \item \textbf{HSV:} Oral acyclovir 400 mg tid or valacyclovir 1 g BID for 7 days; suppressive therapy for recurrent cases
  \item \textbf{Candidiasis:} Clotrimazole troches or nystatin suspension; fluconazole 200 mg daily for refractory cases
  \item \textbf{Chemotherapy mucositis:} Oral cryotherapy (ice chips) during infusion, palifermin (KGF), low-level laser therapy
  \item \textbf{Radiation mucositis:} Benzydamine mouth rinse, topical anesthetics, surcalfate
  \item \textbf{Autoimmune:} High-dose systemic corticosteroids with steroid-sparing immunomodulators (mycophenolate, rituximab)
  \item \textbf{Pain control:} Multimodal approach with topical, oral, and IV analgesics including patient-controlled analgesia for severe mucositis
\end{itemize}

\clearpage
